\chapter{A Linear-Gaussian Generative Realization}
\label{ch:generative-gaussian}

This chapter realizes the fixed directed law of \Cref{ch:generative} in two
associated bundles induced from the same principal $G$-bundle, using possibly
different representations and separately chosen local frames. The Gaussian
is a tractable subclass, not the ambient theory. The specialization assumes that
at a root the model-to-belief kernel induces a map from model laws to belief laws that may realize
$\widetilde\Phi:\mathcal E_m\to\mathcal E_b$.  The sample matrix inside that kernel is not itself
the associated-bundle morphism.  At a nonroot the kernel is parent-indexed unless its domain is
enlarged.  The existence of $\Phi:\mathcal E_b\to\mathcal E_m$ creates no reverse conditional.

\section{Directed forest and conditional laws}
\label{sec:gen-gaussian}

Take
\begin{equation}
k_i\in\R^K,
\qquad
m_i\in\R^{d_m},
\qquad
o_i\in\R^{d_{o,i}},
\label{eq:gen-lg-domains}
\end{equation}
and specialize the directed acyclic graph to a forest, so each nonroot $i$ has one parent
$j=\operatorname{par}(i)$.  This is a displayed subclass; the general DAG replaces each parent
term by a finite sum. \status{DEFINITION}\label{def:gen-lg-forest}

Fix one design point \(c_a\). Let
\(\Theta_{ij}^b:=\Theta_{a,ij}^b\in G\) and
\(\Theta_{ij}^m:=\Theta_{a,ij}^m\in G\) abbreviate the design-indexed
same-channel graph-link copies of \Cref{def:gen-product-action}. A submodel
that constrains one common link across all design points retains the
stabilizer of that constraint, namely the frame rechoices for which
\(h_{a,i}^x\Theta_{ij}^x(h_{a,j}^x)^{-1}\) is independent of \(a\).
Context-independent rechoices form a sufficient, but not generally maximal,
subgroup; any larger law must be declared explicitly. For a nonroot,
\begin{align}
m_i\given m_j,X
&\sim\mathcal N\left(
F_i^m\rho_m(\Theta_{ij}^m)m_j+\eta_i^m,
\Sigma_i^m\right),
\label{eq:gen-lg-model}\\
k_i\given k_j,m_i,X
&\sim\mathcal N\left(
F_i^b\rho_b(\Theta_{ij}^b)k_j+B_im_i+\eta_i^b,
\Sigma_i^b\right),
\label{eq:gen-lg-state}\\
o_i\given k_i,m_i,X
&\sim\mathcal N\left(
H_i^bk_i+H_i^mm_i+\eta_i^o,
\Sigma_i^o\right).
\label{eq:gen-lg-obs}
\end{align}
The parameter shapes are
\begin{align}
&\rho_b(\Theta_{ij}^b),F_i^b\in\R^{K\times K},
&&\rho_m(\Theta_{ij}^m),F_i^m\in\R^{d_m\times d_m},
&&B_i\in\R^{K\times d_m},
\label{eq:gen-lg-shapes-latent}\\
&H_i^b\in\R^{d_{o,i}\times K},
&&H_i^m\in\R^{d_{o,i}\times d_m},
&&\eta_i^b\in\R^K,
\quad
\eta_i^m\in\R^{d_m},
\quad
\eta_i^o\in\R^{d_{o,i}}.
\label{eq:gen-lg-shapes-obs}
\end{align}
The represented links are invertible.  The conditional endomorphisms $F_i^b,F_i^m$ need not be:
positive definiteness of the receiving covariance, not invertibility of the mean map, normalizes a
Gaussian kernel.

The bridge matrix $B_i:\R^{d_m}\to\R^K$ acts on sample coordinates.  It is one parameter of a
Gaussian kernel, not a map between probability-law fibers.  At a nonroot, the kernel also depends
on the parent belief coordinate.  Only after fixing that parent, or enlarging the domain to include
it, does integration against a model law define a belief law.  The left-right covariance law below
is a parameter transformation that makes this induced law map covariant; it does not identify
$B_i$ with $\widetilde\Phi$.  No reading supplies a belief-to-model factor.

Each root $r$ has the proper laws
\begin{align}
m_r\given X
&\sim\mathcal N(\bar m_r,\Sigma_{r,0}^m),
\label{eq:gen-lg-root-model}\\
k_r\given m_r,X
&\sim\mathcal N(B_rm_r+\eta_{r,0}^b,\Sigma_{r,0}^b),
\label{eq:gen-lg-root-state}
\end{align}
followed by \eqref{eq:gen-lg-obs}.  Every receiving covariance is positive definite:
\begin{equation}
\Sigma_i^m,\Sigma_{r,0}^m\in\Sym_{++}^{d_m},
\qquad
\Sigma_i^b,\Sigma_{r,0}^b\in\Sym_{++}^{K},
\qquad
\Sigma_i^o\in\Sym_{++}^{d_{o,i}}.
\label{eq:gen-lg-spd}
\end{equation}
\status{HYPOTHESIS}\label{hyp:gen-lg-spd}

\propositionheading{Normalization and root moments}{prop:gen-lg-normalization}
Under \eqref{eq:gen-lg-spd}, every kernel above is a nondegenerate normalized density and the
assembled directed law is a probability measure.  The root belief coordinate has
\begin{equation}
\bar k_r=B_r\bar m_r+\eta_{r,0}^b,
\qquad
\Sigma_{r,0}^{b,\mathrm{marg}}
=\Sigma_{r,0}^b+B_r\Sigma_{r,0}^mB_r^\top.
\label{eq:gen-lg-root-marginal}
\end{equation}
\status{ESTABLISHED}

More generally, if a root model law is
$s=\mathcal N(\mu_m,C_m)$, then the law map induced by the root kernel is
\begin{equation}
\widetilde\phi_r(s)
=\mathcal N\left(
B_r\mu_m+\eta_{r,0}^b,\quad
\Sigma_{r,0}^b+B_rC_mB_r^\top
\right).
\label{eq:gen-lg-induced-law-map}
\end{equation}
This map, rather than $B_r$ alone, is the candidate local realization of
$\widetilde\Phi$. \status{ESTABLISHED}

\noindent\emph{Proof.}
A Gaussian density with positive definite covariance integrates to one for every mean, so the
normalization proposition of \Cref{ch:generative} applies.  The tower property gives the first
moment and the law of total covariance gives the second identity.  Its added term is positive
semidefinite, so the marginal covariance dominates the conditional covariance. \hfill$\square$

The Gaussian observation kernel may be replaced by any normalized Markov kernel from
$\R^K\times\R^{d_m}$ to the observation space.  That changes the likelihood but not the exact ELBO
identity, which uses only the fixed normalized joint. \status{ESTABLISHED}

\section{Independent channel-frame covariance}
\label{sec:gen-gaussian-covariance}

For the global coordinate formulas in this section, invoke
\Cref{hyp:geo-common-trivializations} on a declared region
\(\mathcal C_0\) containing the contexts under study. Choose
\(\sigma_0:\mathcal C_0\to P|_{\mathcal C_0}\) and restrict each local frame
to \(\mathcal C_i\cap\mathcal C_0\). Then write the independent frame
coordinates as
\begin{equation}
U_i^b:\mathcal C_i\cap\mathcal C_0\to G,
\qquad
U_i^m:\mathcal C_i\cap\mathcal C_0\to G,
\qquad
u_i^x=\sigma_0\cdot(U_i^x)^{-1}.
\label{eq:gen-frames}
\end{equation}
Thus $\sigma_0$ is common only on the explicitly trivialized region; no
global triviality of \(P\) is inferred. The corresponding
relative principal-frame field is
\begin{equation}
h_i=U_i^b(U_i^m)^{-1},
\qquad
u_i^m=u_i^b h_i.
\label{eq:gen-relative-principal-frame}
\end{equation}
This $G$-valued field exists even when $\rho_b$ and $\rho_m$ are inequivalent
or have different dimensions. Its represented actions $\rho_b(h_i)$ and
$\rho_m(h_i)$ live in different endomorphism spaces and do not define the
bridge $B_i$ or either associated-bundle morphism.
An independent pair of channel-frame re-expressions is
$(g_i^b,g_i^m)\in G\times G$, with represented matrices
$R_i^b=\rho_b(g_i^b)$ and $R_i^m=\rho_m(g_i^m)$ as in
\eqref{eq:gen-rep-actions}. The pair records two choices of local
trivialization of bundles associated to the same $P$; it does not replace the
principal structure group $G$ by $G\times G$.
In the convention of \eqref{eq:geo-local-reframing}, the corresponding
principal-section rechoices are
$u_i^{b\prime}=u_i^b\cdot(g_i^b)^{-1}$ and
$u_i^{m\prime}=u_i^m\cdot(g_i^m)^{-1}$. This accounts for the opposite-sided
inverse in the geometric local-map law and makes the formulas below identical
to \eqref{eq:geo-tildephi-gauge-law}.

The same-channel links obey \eqref{eq:gen-gauge-links}.  The conditional maps and offsets transform
as
\begin{align}
F_i^{b\prime}
&=R_i^bF_i^b(R_i^b)^{-1},
&
F_i^{m\prime}
&=R_i^mF_i^m(R_i^m)^{-1},
\label{eq:gen-gauge-selfmaps}\\
B_i'
&=R_i^bB_i(R_i^m)^{-1},
&
\eta_i^{b\prime}
&=R_i^b\eta_i^b,
\qquad
\eta_i^{m\prime}=R_i^m\eta_i^m.
\label{eq:gen-gauge-bridge-offsets}
\end{align}
The left-right law for $B_i$, together with the covariance and offset laws, makes the induced root
law map \eqref{eq:gen-lg-induced-law-map} obey the local law of $\widetilde\Phi$.  The parameter
$B_i$ alone is not that morphism, and the two frames remain independent.  Receiving covariances
transform by their own congruences:
\begin{equation}
\Sigma_i^{b\prime}=R_i^b\Sigma_i^b(R_i^b)^\top,
\qquad
\Sigma_i^{m\prime}=R_i^m\Sigma_i^m(R_i^m)^\top.
\label{eq:gen-gauge-covariances}
\end{equation}
For either channel,
\begin{equation}
\Lambda_i=\Sigma_i^{-1},
\qquad
\Sigma_i'=R_i\Sigma_iR_i^\top,
\qquad
\Lambda_i'=R_i^{-\top}\Lambda_iR_i^{-1}.
\label{eq:gen-gauge-precision}
\end{equation}
At a root,
\begin{align}
\bar m_r'
&=R_r^m\bar m_r,
&
\Sigma_{r,0}^{m\prime}
&=R_r^m\Sigma_{r,0}^m(R_r^m)^\top,
\label{eq:gen-gauge-root-model}\\
B_r'
&=R_r^bB_r(R_r^m)^{-1},
&
\eta_{r,0}^{b\prime}
&=R_r^b\eta_{r,0}^b,
\qquad
\Sigma_{r,0}^{b\prime}
=R_r^b\Sigma_{r,0}^b(R_r^b)^\top.
\label{eq:gen-gauge-root-state}
\end{align}
Gauge-inert observations require
\begin{equation}
H_i^{b\prime}=H_i^b(R_i^b)^{-1},
\qquad
H_i^{m\prime}=H_i^m(R_i^m)^{-1},
\qquad
\eta_i^{o\prime}=\eta_i^o,
\qquad
\Sigma_i^{o\prime}=\Sigma_i^o.
\label{eq:gen-gauge-observation}
\end{equation}
\status{DEFINITION}

\propositionheading{Gaussian kernel covariance}{prop:gen-lg-kernel-covariance}
The transformation laws above satisfy the three kernel pushforward identities
\eqref{eq:gen-gauge-pushforward-model}--\eqref{eq:gen-gauge-pushforward-obs}, and hence the
normalized joint and its evidence obey
\eqref{eq:gen-gauge-joint-pushforward}--\eqref{eq:gen-gauge-evidence-invariance} under all
independent channel-frame re-expressions. \status{ESTABLISHED}

\noindent\emph{Proof.}
For the belief kernel, the transformed conditional mean evaluated at transformed inputs is
\begin{equation}
F_i^{b\prime}\rho_b(\Theta_{ij}^{b\prime})(R_j^bk_j)
+B_i'(R_i^mm_i)+\eta_i^{b\prime}
=R_i^b\left(
F_i^b\rho_b(\Theta_{ij}^b)k_j+B_im_i+\eta_i^b
\right).
\label{eq:gen-gauge-mean-identity}
\end{equation}
The model-channel calculation is the same without a bridge.  Congruence
\eqref{eq:gen-gauge-covariances} therefore makes each transformed Gaussian the receiving-frame
pushforward of the original.  At the observation kernel,
\begin{equation}
H_i^{b\prime}(R_i^bk_i)+H_i^{m\prime}(R_i^mm_i)+\eta_i^{o\prime}
=H_i^bk_i+H_i^mm_i+\eta_i^o,
\label{eq:gen-gauge-observation-identity}
\end{equation}
and its covariance is fixed.  The root calculation is identical.  Composition now invokes the
frame-covariance result of \Cref{ch:generative}. \hfill$\square$

Equation~\eqref{eq:gen-lg-root-marginal} transforms by belief-channel congruence, because both of
its summands do.  Positive definiteness is preserved independently in each channel.

Group-valued frame coordinates remain primary.  An expression $U=\exp\phi$ is valid only inside an
exponential chart; for example,
$\operatorname{diag}(-1,-2)\in\GL^+(2)$ has no real logarithm because each negative eigenvalue
occurs in one Jordan block \citep{culver1966existence}. \status{ESTABLISHED}

\section{Common-frame specialization}
\label{sec:gen-gaussian-diagonal}

Under the common-frame specialization of \Cref{sec:geo-diagonal}, impose
$u_i^b=u_i^m$ and use one frame re-expression $g_i\in G$ with
\begin{equation}
R_i^b=\rho_b(g_i),
\qquad
R_i^m=\rho_m(g_i).
\label{eq:gen-gaussian-diagonal-action}
\end{equation}
If the graph links are also identified, $\Theta_{ij}^b$ and
$\Theta_{ij}^m$ are the same group link seen through the two
representations. Equations
\eqref{eq:gen-gauge-selfmaps}--\eqref{eq:gen-gauge-observation} then reduce to the familiar
common-frame, two-representation formulas. \status{ESTABLISHED}

This reduction ties the gauge choices but does not identify the latent dimensions, make $B_i$
invertible, turn $\Phi$ and $\widetilde\Phi$ into inverses, or add a reverse generative factor.
Compatibility of a cross map with the two connections remains the parallelness
condition of \Cref{ch:geometry}; an equal connection is a further hypothesis.

Finally, the assembled Gaussian precision need not have the Laplacian-plus-self-terms form used
later.  A general directed transition can contribute an off-diagonal block that is neither
symmetric nor sign-correct.  The interaction family is therefore a separately declared Gaussian
operator subfamily, not a consequence of linear-Gaussian directedness. \status{NOT-CLAIMED}
